\section{Малые размышления}

\subsection{Демократия невообразимого}
С человеческой точки зрения гугол, гуголплекс, уровень \(d\), уровень \(p\), крипто кома и
даже число Грэма сливаются в одну размытую категорию: «примерно бесконечность». У интуиции
есть потолок --- когнитивный горизонт, --- и лежит он абсурдно низко, где-то в районе
нескольких тысяч. За ним все великаны на одно лицо. Разница между ними огромна, реальна и
нам невидима.

\subsection{Имя вместо числа}
Крипто кому нельзя удержать --- поэтому мы держим её имя: \(\ccoma\), три штриха. Имя
конечно, число не записать. Назвать --- способ приручить то, что не унести. Но поводок ---
это не зверь. Что у нас в руках, когда мы пишем \(\ccoma\), --- число или слово, стоящее
там, куда число не помещается?

\subsection{Почти все великаны безымянны}
Коротких описаний конечное число: фраз короче заданной длины попросту не бесконечно много.
А чисел --- бесконечно. Значит, почти всякое число, включая почти всякое число размера
\(d\), нельзя ни назвать, ни описать, ни выделить. Крипто кома --- редкий именованный
великан в бесконечной толпе анонимных.

\subsection{Физически трансцендентное}
Компьютер, работающий до тепловой смерти Вселенной и переключающий бит каждый планковский
миг, не напечатал бы цифры уровня \(d\) --- даже число его цифр. У вычислений есть физический
предел. Значит, некоторые конечные числа реальны, но навсегда неосуществимы: истинны --- и
никогда не будут построены, сосчитаны или показаны нигде в этом космосе.

\subsection{Лестница без вершины}
\(z\) --- последняя буква, а не последнее число. Примените \(\star\) ещё раз --- и вы прошли
крипто кому; ещё раз --- и она осталась далеко внизу. Назовите любой потолок, и вы только что
построили пол следующей комнаты. Наибольшего числа нет --- есть лишь самое большое из тех,
что вы удосужились назвать.

\subsection{Открыли или придумали?}
Ждала ли крипто кома в платоновском небе двух мальчиков на лавочке? Или вспыхнула в бытие в
тот миг, когда её записали? Ультрафинитисты сказали бы: чисел, до которых нам не добраться,
попросту нет. Мальчики весело изобрели кошмар для философии.
